\section{Results}
% Budget: 2.3 pages, including all main tables and figures.
% Order: data completeness; sampling and accuracy; feature agreement;
% correction; quality gate; synthetic consequences and validation diagnostics.
% Define estimable, adjusted decision, correction support, and unavailable estimate.
% Do not insert numbers until outputs exist. Include every planned result.
% Claim rules: bench protocol section 16 and consequence protocol section 11.
% A feature claim needs all prespecified programs and the Holm-adjusted decision.
% Camera references apply only to duration, path length, mean speed, and loop height.
% Speed CV and band power concern commanded movement within the executed-motion check.
% Keep per-program results descriptive. A failed test does not prove a difference.
% Correction residuals are a lack-of-fit check, not independent transfer validation.
% Separate grid assumptions from measured-input consequences throughout.
% Do not claim human agreement, reliability, screening accuracy, or clinical utility.

\subsection{Completeness, sampling, and accuracy}
\pending{R-completeness: report valid runs, failures, reruns, missing camera records, deviations, and non-estimable programs.}
\begin{table}[t]
% Caption specification: State the observed sampling and accuracy pattern within the tested range.
% Show device and program summaries for sample intervals, gap shares, noise floors, camera distances, and command trace error.
% Give units, valid-run counts, and ranges.
\caption{}
\label{tab:sampling}
\pending{T-sampling: insert the sampling and spatial accuracy table. Write the caption from its specification.}
\end{table}
\pending{R-sampling: describe measured sampling and accuracy without an acceptance claim. Include hold sample counts and camera validity.}

\subsection{Feature agreement}
\begin{table}[t]
% Caption specification: State which features meet the prespecified equivalence decision across their full program sets.
% Show limits, comparison scales, camera-reference status, feature p-values, adjusted and unadjusted decisions, and estimability.
% Keep per-program decisions descriptive.
\caption{}
\label{tab:agreement}
\pending{T-agreement: insert all feature decisions, including failures. Write the caption from its specification.}
\end{table}
\begin{figure}[t]
\pending{F-agreement: insert the bias intervals and repeated-measures agreement plots. Write the caption from its specification.}
% Caption specification: State the observed pattern of bias relative to the prespecified limits.
% Show program-specific intervals and Bland--Altman points with mounting-mean and single-run limits of agreement.
% Identify comparison scales and the camera-referenced features.
\caption{}
\label{fig:agreement}
\end{figure}
\pending{R-agreement: state only supported feature-level decisions. Report per-program bias and sensitivity results as descriptive evidence.}

\subsection{Correction and quality gate}
\begin{table}[t]
% Caption specification: State whether each correction passes its lack-of-fit check across every prespecified program.
% Show model terms, fitted or fixed status, coefficient uncertainty, and omitted-mounting residual intervals against the same limits.
% Do not describe these intervals as independent validation.
\caption{}
\label{tab:correction}
\pending{T-correction: insert coefficients, residuals, and support status. Write the caption from its specification.}
\end{table}
\pending{R-correction: report residual sizes, failed checks, and unavailable fits. Bound use to the tested measurement chain and range.}
\begin{table}[t]
% Caption specification: State which designed faults the quality gate detects on each device.
% Show main-run acceptance and detection by fault and device, with exact intervals and denominators.
% State the locked Pencil dwell expectation.
\caption{}
\label{tab:gate}
\pending{T-gate: insert acceptance and fault-detection summaries. Write the caption from its specification.}
\end{table}
\pending{R-gate: describe each fault outcome. Do not infer detection of untested faults.}

\subsection{Synthetic consequences}
\begin{table}[t]
% Caption specification: State the scenario-specific primary changes from the paired synthetic reference.
% Show raw change, corrected change, absolute reduction, uncertainty, material-change share, usable pairs, and correction support for every feature.
% Primary summaries: participant-wise AUC for all-W against all-P; W-to-P calibration
% intercept for transfer; participant-wise AUC under the null effect for confounding.
% Report D_raw and absolute reduction, as defined in Methods, for each summary.
% Separate assumed grids from measured inputs.
\caption{}
\label{tab:consequence}
\pending{T-consequence: insert the primary consequence summaries and write the caption from its specification.}
\end{table}
\begin{figure}[t]
\pending{F-consequence: insert bias-response, transfer-calibration, and decision-curve panels. Write the caption from its specification.}
% Caption specification: State how the measurement-chain difference changes synthetic model metrics under the specified assumptions.
% Show uncorrected and corrected bias response, directional transfer calibration, and decision curves against paired references.
% Distinguish mean uncertainty from simulation distribution intervals. Retain unavailable levels and corrections that worsen results.
\caption{}
\label{fig:consequence}
\end{figure}
% Report the target geometry, scale, speed, frequency, and selected program for every feature.
% Put scale-and-speed distance beside absolute frequency difference. State range status.
% Mark extrapolation as sensitivity evidence. Do not select programs by performance.
\pending{R-program-selection: report the selected programs, target settings, both distances, and range status.}
\pending{R-consequence: report the scenario-specific primary metrics and all supporting metrics. Explain negative absolute reductions and incomplete cells.}
\pending{R-validation: compare record-wise, participant-wise, and both device-held-out directions. Report overlap, baselines, refusals, and sensitivity analyses.}
% State beside the affected results that path length, speed CV, and band-power population
% scaling has no independent population justification. Joint features are not one drawing.
% A small synthetic effect does not establish equivalence. No numerical lower bound
% on human AUC or calibration distortion follows from omitted surface behavior.
% Optional example: omit by default. Add a separate result only after written clearance,
% public provenance, participant grouping, deduplication, and full citations are verified.
